\documentclass{report}
\usepackage{amsmath,amssymb}
\setlength{\parindent}{0mm}
\setlength{\parskip}{1em}
\begin{document}
\begin{center}
\rule{6in}{1pt} \
{\large Heidi Thornquist \\
{\bf Enabling Next-Generation Circuit Simulation Using Trilinos}}

Sandia National Laboratories \\ PO Box 5800 \\ Albuquerque \\ NM 87185
\\
{\tt hkthorn@sandia.gov}\\
Erik Boman\\
Eric Keiter\\
Siva Rajamanickam\\
Rich Schiek\end{center}

The Xyce Parallel Circuit Simulator, which has demonstrated scalable
circuit simulation on hundreds of processors, heavily leverages the
high-performance scientific libraries provided by Trilinos. With the move
towards multi-core CPUs and GPU technology, retaining this scalability on
future parallel architectures will be a challenge. The computational
expense in traditional analog circuit simulation is in repeatedly solving
linear systems of equations, which are at the center of a nested solver
loop. Solving these linear systems requires their assembly, which depends
upon device evaluations for the whole circuit. So, the computational
expense is dominated by either the device evaluations or the numerical
method used to solve the linear systems. This talk will present the
enabling technologies provided by Trilinos to exploit both coarse and
fine-grained parallelism, including a new effort in hybrid linear
solvers, that will allow Xyce to retain scalable performance on
next-generation architectures.


\end{document}
